\documentclass[11pt]{article}
\usepackage[a4paper,margin=2.4cm]{geometry}
\usepackage[T1]{fontenc}
\usepackage{lmodern}
\usepackage{amsmath,amssymb}
\newcommand{\dphi}[1]{\Delta\phi_{#1}}
\newcommand{\half}{\tfrac{1}{2}}
\title{MOS Capacitor Simulator\\[4pt]\large Equations as implemented in the code}
\date{}
\begin{document}
\maketitle

\section{Model and notation}
The structure is gate -- oxide -- silicon (p-type or n-type) -- grounded bulk contact, in one dimension $x$.
\begin{align}
V_t &= \frac{kT}{q}, & \phi &= \frac{V}{V_t}, & E &= -\frac{dV}{dx}
\end{align}
The grid nodes are $x_0\dots x_N$. Node $0$ is the gate, node $s$ is the oxide--semiconductor interface and
node $N$ is the bulk contact. The oxide mesh is uniform; in the semiconductor the spacing grows geometrically:
\begin{align}
\Delta x_i &= x_{i+1}-x_i, & \Delta x_{s+k} &= a\,b^{k}
\quad\Longrightarrow\quad x_{s+k} = t_{ox} + a\,\frac{b^{k}-1}{b-1}
\end{align}
Each node owns a control volume. Its semiconductor part has width
\begin{align}
h_i = \begin{cases} 0 & i<s \quad\text{(oxide)}\\[2pt]
\half\Delta x_s & i=s \quad\text{(interface: half a cell)}\\[2pt]
\half\left(\Delta x_{i-1}+\Delta x_i\right) & s<i<N \end{cases}
\end{align}
Differences of the normalised potential between neighbouring nodes are written
$\dphi{i}=\phi_{i+1}-\phi_i$, and $n_{int}$ is the intrinsic carrier concentration.

\section{Current densities (Scharfetter--Gummel)}
\begin{align}
J_n &= q\mu_n nE + qD_n\frac{dn}{dx} = kT\mu_n\left(-n\frac{d\phi}{dx}+\frac{dn}{dx}\right)\\
J_p &= q\mu_p pE - qD_p\frac{dp}{dx} = -kT\mu_p\left(p\frac{d\phi}{dx}+\frac{dp}{dx}\right)
\end{align}
The current is a flux between two nodes, so it is defined on the element $i+\half$ between nodes $i$ and
$i+1$. Taking the potential linear inside an element and writing $n=u\,e^{\phi}$ gives
$J_n = kT\mu_n e^{\phi}\,du/dx$. Integrating $J_n e^{-\phi}$ over the element with constant $J_n$:
\begin{align}
J_{n,i+\half} &= \frac{kT\mu_n}{\Delta x_i}\Big[B(\dphi{i})\,n_{i+1}-B(-\dphi{i})\,n_i\Big] \label{eq:jn}
\end{align}
with the Bernoulli function
\begin{align}
B(x)=\frac{x}{e^{x}-1}, \qquad B(x)\,e^{x}=B(-x)
\end{align}
For holes take $p=u\,e^{-\phi}$, so $J_p=-kT\mu_p e^{-\phi}\,du/dx$ and
\begin{align}
J_{p,i+\half} &= \frac{kT\mu_p}{\Delta x_i}\Big[B(\dphi{i})\,p_i-B(-\dphi{i})\,p_{i+1}\Big] \label{eq:jp}
\end{align}

\subsection{Quasi-Fermi form}
With the quasi-Fermi potentials $\Phi_n$, $\Phi_p$ defined by
\begin{align}
n=n_{int}\exp\!\left(\frac{V-\Phi_n}{V_t}\right),\qquad p=n_{int}\exp\!\left(\frac{\Phi_p-V}{V_t}\right) \label{eq:qf}
\end{align}
the currents \eqref{eq:jn}, \eqref{eq:jp} become, exactly,
\begin{align}
\frac{J_{n,i+\half}}{q}&=-\frac{V_t\mu_n}{\Delta x_i}\,B(\dphi{i})\,n_{i+1}\left(e^{d_n}-1\right),
& d_n&=\frac{\Phi_{n,i+1}-\Phi_{n,i}}{V_t} \label{eq:jnqf}\\
\frac{J_{p,i+\half}}{q}&=-\frac{V_t\mu_p}{\Delta x_i}\,B(\dphi{i})\,p_i\left(e^{d_p}-1\right),
& d_p&=\frac{\Phi_{p,i+1}-\Phi_{p,i}}{V_t} \label{eq:jpqf}
\end{align}
In this form drift and diffusion are never subtracted from each other, so a small net current is computed
accurately. A current is zero exactly when the quasi-Fermi potential is flat. This form is used in the
small-signal and transient solvers.

\section{Steady state (equilibrium and DC bias)}
\subsection{Poisson's equation}
\begin{align}
\frac{d}{dx}\left(\varepsilon\frac{dV}{dx}\right) = -\rho, \qquad \rho = q\left(N_D-N_A+p-n\right)
\end{align}
Integrating over the control volume of node $i$ (there is no charge in the oxide, where $h_i=0$):
\begin{align}
F_i \equiv \varepsilon_{i+\half}\frac{V_{i+1}-V_i}{\Delta x_i}-\varepsilon_{i-\half}\frac{V_i-V_{i-1}}{\Delta x_{i-1}}
+q\,h_i\left(N_D-N_A+p_i-n_i\right)+Q_f\,\delta_{i,s}=0 \label{eq:poisson}
\end{align}
where $\varepsilon_{i+\half}=\varepsilon_{ox}$ for $i<s$ and $\varepsilon_{si}$ for $i\ge s$. At the interface
node $i=s$ the two permittivities differ, which enforces continuity of $\varepsilon E$; the fixed oxide
charge $Q_f$ (per unit area) sits at this node only.

Boundary conditions. $V$ is measured from the bulk, so $V_N=0$. The gate and the semiconductor have
different work functions, so a gate voltage $V_G$ puts the gate node at
\begin{align}
V_0=V_G-\phi_{ms},\qquad \phi_{ms}=\phi_m-\phi_s,\qquad
\phi_s=\chi+\frac{E_g}{2q}-V_t\ln\frac{n_{bulk}}{n_{int}}
\end{align}
with $\phi_m$ the gate work function, $\chi$ the electron affinity and $n_{bulk}$ the electron
concentration at the bulk contact. This covers both substrate types: $N_A>N_D$ is p-type
($n_{bulk}<n_{int}$) and $N_D>N_A$ is n-type. The bands are flat ($V=0$ everywhere) at the flat-band voltage
\begin{align}
V_{FB}=\phi_{ms}-\frac{Q_f}{C_{ox}},\qquad C_{ox}=\frac{\varepsilon_{ox}}{t_{ox}}
\end{align}
Equilibrium is $V_G=0$, which is in general not the flat-band condition.

The equation is non-linear because $n$ and $p$ depend exponentially on $V$. One Newton step
$J\,\Delta V=-F$ is taken with the carriers following the potential at fixed quasi-Fermi potentials:
\begin{align}
J_{i,i+1} = \frac{\varepsilon_{i+\half}}{\Delta x_{i}},\qquad J_{i,i-1} = \frac{\varepsilon_{i-\half}}{\Delta x_{i-1}},\qquad
J_{i,i} = -\frac{\varepsilon_{i+\half}}{\Delta x_i}-\frac{\varepsilon_{i-\half}}{\Delta x_{i-1}}-\frac{q\,h_i\,(n_i+p_i)}{V_t}
\end{align}
followed by $V\leftarrow V+\Delta V$, $n\leftarrow n\,e^{\Delta V/V_t}$, $p\leftarrow p\,e^{-\Delta V/V_t}$.
Steps larger than $V_t$ are damped to $\pm V_t\ln\!\left(1+|\Delta V|/V_t\right)$.

\subsection{Carrier continuity}
\begin{align}
\frac{\partial n}{\partial t}=\frac1q\frac{\partial J_n}{\partial x}-R,\qquad
\frac{\partial p}{\partial t}=-\frac1q\frac{\partial J_p}{\partial x}-R \label{eq:continuity}
\end{align}
Recombination is Shockley--Read--Hall:
\begin{align}
R=\frac{np-n_{int}^2}{S},\qquad S=\tau_p\,(n+n_{int})+\tau_n\,(p+n_{int}) \label{eq:srh}
\end{align}
In steady state, integrating over the control volume of node $i$:
\begin{align}
\frac1q\left(J_{n,i+\half}-J_{n,i-\half}\right)=h_iR_i,\qquad
-\frac1q\left(J_{p,i+\half}-J_{p,i-\half}\right)=h_iR_i
\end{align}
Inserting \eqref{eq:jn}, \eqref{eq:jp} and dividing by $V_t\mu$, with $S_i$ and the other carrier taken
from the previous iteration:
\begin{align}
\alpha_{i-1}n_{i-1}+\alpha_in_i+\alpha_{i+1}n_{i+1} &= -\frac{h_i\,n_{int}^2}{V_t\mu_nS_i}\\
\alpha_{i+1}=\frac{B(\dphi{i})}{\Delta x_i},\quad
\alpha_{i-1}&=\frac{B(-\dphi{i-1})}{\Delta x_{i-1}},\quad
\alpha_i=-\left[\frac{B(-\dphi{i})}{\Delta x_i}+\frac{B(\dphi{i-1})}{\Delta x_{i-1}}\right]-\frac{h_i\,p_i}{V_t\mu_nS_i}
\end{align}
\begin{align}
\beta_{i-1}p_{i-1}+\beta_ip_i+\beta_{i+1}p_{i+1} &= -\frac{h_i\,n_{int}^2}{V_t\mu_pS_i}\\
\beta_{i+1}=\frac{B(-\dphi{i})}{\Delta x_i},\quad
\beta_{i-1}&=\frac{B(\dphi{i-1})}{\Delta x_{i-1}},\quad
\beta_i=-\left[\frac{B(\dphi{i})}{\Delta x_i}+\frac{B(-\dphi{i-1})}{\Delta x_{i-1}}\right]-\frac{h_i\,n_i}{V_t\mu_pS_i}
\end{align}
Boundary conditions:
\begin{itemize}
\item Interface ($i=s$): no current enters the oxide, $J_{s-\half}=0$. The equations above are used with the
$i-1$ terms dropped and $h_s=\half\Delta x_s$.
\item Bulk contact ($i=N$): $n_N$, $p_N$ keep the charge-neutral equilibrium values, from
$p-n=N_A-N_D$ and $np=n_{int}^2$.
\end{itemize}
Poisson's equation and the two continuity equations are solved one after the other and repeated until
$V$, $n$ and $p$ stop changing. The gate voltage is moved in steps of at most $0.1$~V, each step starting
from the previous solution.

\subsection{Solving the continuity equations accurately}\label{sec:flux}
Both equations have the form
$\mathcal F_i-\mathcal F_{i-1}-\sigma_iy_i=-g_i$ with flux $\mathcal F_i=u_iy_{i+1}-d_iy_i$ and
$u_i,d_i,\sigma_i,g_i\ge0$. The fluxes are up to $10^{17}$ times larger than the recombination terms, so
ordinary Gaussian elimination loses them to round-off. The elimination is therefore arranged so that it
never subtracts: after eliminating nodes $s\dots i$, node $i$ obeys $u_iy_{i+1}-(d_i+e_i)\,y_i=-b_i$ with
\begin{align}
e_s=\sigma_s,\quad b_s=g_s,\qquad
e_{i+1}=\sigma_{i+1}+\frac{u_i\,e_i}{d_i+e_i},\qquad b_{i+1}=g_{i+1}+\frac{d_i\,b_i}{d_i+e_i}
\end{align}
and the back substitution is $y_i=\left(u_iy_{i+1}+b_i\right)/\left(d_i+e_i\right)$.

\subsection{Deep depletion}
The steady state always contains the inversion layer. If the gate is swept from accumulation
towards inversion faster than minority carriers can be generated, no inversion layer forms and the
depletion region keeps widening. This fast-sweep limit is modelled by freezing the minority carrier:
\begin{itemize}
\item the sweep starts from the steady state at the accumulation end of the gate-voltage range;
\item the minority carrier (electrons for p-type, holes for n-type) keeps the profile it has there. It is
not solved from the continuity equation and not rescaled in the Newton step, so it also drops out of the
Jacobian: $J_{i,i}$ contains only the majority carrier;
\item the majority carrier and Poisson's equation are solved as before.
\end{itemize}
The capacitance equals the high-frequency one up to threshold and keeps falling beyond it, as
$C=\left(1/C_{ox}+W/\varepsilon_{si}\right)^{-1}$ with $W$ the depletion width.

\section{Small-signal AC analysis}\label{sec:ac}
A small sinusoidal signal is added to the DC gate bias, and every quantity is written as its DC value
(subscript $0$) plus a phasor:
\begin{align}
V_G&=V_{DC}+\tilde v\,e^{j\omega t}, & V&=V_0+\tilde V e^{j\omega t}, &
\Phi_n&=\Phi_{n0}+\tilde W_n e^{j\omega t}, & \Phi_p&=\Phi_{p0}+\tilde W_p e^{j\omega t}
\end{align}
Only terms linear in the perturbation are kept, and $\partial/\partial t\to j\omega$. From \eqref{eq:qf}
the carrier phasors are
\begin{align}
\tilde n=\frac{n_0}{V_t}\left(\tilde V-\tilde W_n\right),\qquad
\tilde p=\frac{p_0}{V_t}\left(\tilde W_p-\tilde V\right)
\end{align}

\subsection{Linearised equations}
No DC current flows in a MOS capacitor, so the DC quasi-Fermi potentials are flat ($d_n=d_p=0$) and
$B(\dphi{0,i})\,n_{0,i+1}=B(-\dphi{0,i})\,n_{0,i}\equiv\bar n_i$,
$B(\dphi{0,i})\,p_{0,i}=B(-\dphi{0,i})\,p_{0,i+1}\equiv\bar p_i$.
Linearising \eqref{eq:jnqf}, \eqref{eq:jpqf} then leaves plain conductances:
\begin{align}
\frac{\tilde J_{n,i+\half}}{q}=-\frac{\mu_n\bar n_i}{\Delta x_i}\left(\tilde W_{n,i+1}-\tilde W_{n,i}\right),\qquad
\frac{\tilde J_{p,i+\half}}{q}=-\frac{\mu_p\bar p_i}{\Delta x_i}\left(\tilde W_{p,i+1}-\tilde W_{p,i}\right)
\end{align}
Linearising \eqref{eq:srh} around the DC solution ($R_0$, $S_0$ are the DC values):
\begin{align}
\tilde R=r_n\tilde n+r_p\tilde p,\qquad r_n=\frac{p_0-R_0\tau_p}{S_0},\qquad r_p=\frac{n_0-R_0\tau_n}{S_0}
\end{align}
The three equations of node $i$ are the linearised \eqref{eq:poisson} and \eqref{eq:continuity}:
\begin{align}
\varepsilon_{i+\half}\frac{\tilde V_{i+1}-\tilde V_i}{\Delta x_i}-\varepsilon_{i-\half}\frac{\tilde V_i-\tilde V_{i-1}}{\Delta x_{i-1}}
+q\,h_i\left(\tilde p_i-\tilde n_i\right)&=0 \label{eq:poissonac}\\
-\frac1q\left(\tilde J_{n,i+\half}-\tilde J_{n,i-\half}\right)+h_i\left(\tilde R_i+j\omega\,\tilde n_i\right)&=0 \label{eq:nac}\\
-\frac1q\left(\tilde J_{p,i+\half}-\tilde J_{p,i-\half}\right)-h_i\left(\tilde R_i+j\omega\,\tilde p_i\right)&=0 \label{eq:pac}
\end{align}
$\phi_{ms}$ and $Q_f$ are constants, so they do not appear.

\subsection{Boundary conditions}
\begin{itemize}
\item Gate: $\tilde V_0=\tilde v$. The oxide has no charge, so $\tilde V$ is linear in it and the oxide acts as
the capacitor $C_{ox}$ between gate and interface.
\item Interface ($i=s$): $\tilde J_{n,s-\half}=\tilde J_{p,s-\half}=0$, with $h_s=\half\Delta x_s$.
\item Bulk contact ($i=N$): $\tilde V_N=\tilde W_{n,N}=\tilde W_{p,N}=0$.
\end{itemize}

\subsection{Solution}\label{sec:block}
The AC problem is linear, so it is solved once per bias point and frequency, with no iteration. With
$X_i=(\tilde V_i,\tilde W_{n,i},\tilde W_{p,i})^T$, the three equations of node $i$ are
\begin{align}
\big(U_iX_{i+1}-D_iX_i\big)-\big(U_{i-1}X_i-D_{i-1}X_{i-1}\big)-L_iX_i=-g_i
\end{align}
with $3\times3$ complex matrices: $U_i=D_i=\mathrm{diag}\!\left(\varepsilon_{si},\ \mu_n\bar n_i,\ \mu_p\bar p_i\right)/\Delta x_i$
hold the flux coefficients of element $i+\half$, $L_i$ the local terms (charge, recombination, $j\omega$),
and $g_i$ the sources. The oxide enters at the interface node only: $C_{ox}$ is added to the first diagonal
entry of $L_s$, and $g_s=(C_{ox}\tilde v,0,0)^T$; all other $g_i=0$.
This is the block version of Section~\ref{sec:flux}: the fluxes are never subtracted from each
other. After eliminating nodes $s\dots i$, node $i$ obeys $U_iX_{i+1}-(D_i+E_i)X_i=-b_i$ with
\begin{align}
E_s&=L_s, & b_s&=g_s\\
E_{i+1}&=L_{i+1}+E_i\,(D_i+E_i)^{-1}U_i, & b_{i+1}&=g_{i+1}+D_i\,(D_i+E_i)^{-1}b_i
\end{align}
and back substitution $X_i=(D_i+E_i)^{-1}\left(U_iX_{i+1}+b_i\right)$ starting from $X_N=0$.

\subsection{Capacitance and conductance}
The gate charge per unit area is $Q=\varepsilon_{ox}E_{ox}$, so its phasor is
$\tilde Q=C_{ox}(\tilde v-\tilde V_s)$ and
\begin{align}
\tilde C=\frac{\tilde Q}{\tilde v}=C_{ox}\left(1-\frac{\tilde V_s}{\tilde v}\right),\qquad
C=\mathrm{Re}\,\tilde C,\qquad \frac{G}{\omega}=-\mathrm{Im}\,\tilde C
\end{align}
$C$ is the capacitance and $G$ the conductance per unit area at angular frequency $\omega$ (the admittance
is $j\omega\tilde C$). At low frequency the inversion layer follows the signal and $C$ returns to about
$C_{ox}$ in inversion; at high frequency it cannot, and $C$ stays at its minimum. $G/\omega$ peaks at the
frequency where the change-over happens.

Two kinds of sweep are done: a bias sweep ($C$ and $G/\omega$ versus $V_G$ at a few
frequencies) and a frequency sweep ($C$ and $G/\omega$ versus $f$ at a few gate voltages).

\section{Transient simulation}
The steady state and the frozen-carrier deep depletion are the two limits of a gate that moves infinitely
slowly or infinitely fast. In between, the equations are solved in time.

\subsection{Time stepping}
The unknowns at each node are $X_i=(V_i,\Phi_{n,i},\Phi_{p,i})$, with the currents in the form
\eqref{eq:jnqf}, \eqref{eq:jpqf} and $np-n_{int}^2=n_{int}^2\left(e^{(\Phi_p-\Phi_n)/V_t}-1\right)$ in
\eqref{eq:srh}. Backward Euler from time $t$ to $t+\Delta t$, with all terms at the new time except
$n^{old}$, $p^{old}$:
\begin{align}
\varepsilon_{i+\half}\frac{V_{i+1}-V_i}{\Delta x_i}-\varepsilon_{i-\half}\frac{V_i-V_{i-1}}{\Delta x_{i-1}}
+q\,h_i\left(N_D-N_A+p_i-n_i\right)+Q_f\,\delta_{i,s}&=0\\
-\frac1q\left(J_{n,i+\half}-J_{n,i-\half}\right)+h_i\left(R_i+\frac{n_i-n_i^{old}}{\Delta t}\right)&=0\\
-\frac1q\left(J_{p,i+\half}-J_{p,i-\half}\right)-h_i\left(R_i+\frac{p_i-p_i^{old}}{\Delta t}\right)&=0
\end{align}
The boundary conditions are those of the steady state, with the gate node at $V_G(t+\Delta t)-\phi_{ms}$.
The three equations are solved together by Newton's method. The linear system of each Newton iteration
has the block form of Section~\ref{sec:block}, with $U_i$, $D_i$, $L_i$ the derivatives of the fluxes and
local terms with respect to $X$ and the residual as $g_i$, and is solved by the same elimination. Updates
larger than $5V_t$ are clipped.

\subsection{Simulations}
Both start from the steady state at the accumulation end of the gate-voltage range. The capacitance
reported is the small-signal one of Section~\ref{sec:ac} at the highest frequency, around the state at
each time.
\begin{itemize}
\item Sweep-rate dependent C--V: the gate is ramped towards inversion at a rate $dV_G/dt$, with
$\Delta t=\Delta V_G/(dV_G/dt)$ per step. Slow ramps follow the equilibrium high-frequency curve and fast
ramps the deep-depletion curve. The change-over is where the charging current $C_{ox}\,dV_G/dt$ equals the
generation current, about $q\,n_{int}W/(2\tau)$.
\item C--t: the gate is stepped into inversion in a very short time and then held. The device starts in
deep depletion and the capacitance recovers to the equilibrium high-frequency value as generation builds up
the inversion layer. The time steps grow geometrically after the step.
\end{itemize}

\end{document}
